\documentclass[a4paper]{article}
\usepackage{a4wide,amssymb,epsfig,latexsym,multicol,array,hhline,fancyhdr}
\usepackage{vntex}
\usepackage{amsmath}
\usepackage{lastpage}
\usepackage[lined,boxed,commentsnumbered]{algorithm2e}
\usepackage{enumerate}
\usepackage{color}
\usepackage{graphicx}							% Standard graphics package
\usepackage{array}
\usepackage{tabularx, caption}
\usepackage{multirow}
\usepackage{multicol}
\usepackage{rotating}
\usepackage{graphics}
\usepackage{geometry}
\usepackage{setspace}
\usepackage{epsfig}
\usepackage{tikz}
\usetikzlibrary{arrows,snakes,backgrounds}
\usepackage{hyperref}
\hypersetup{urlcolor=blue,linkcolor=black,citecolor=black,colorlinks=true} 
%\usepackage{pstcol} 								% PSTricks with the standard color package

\newtheorem{theorem}{{\bf Theorem}}
\newtheorem{property}{{\bf Property}}
\newtheorem{proposition}{{\bf Proposition}}
\newtheorem{corollary}[proposition]{{\bf Corollary}}
\newtheorem{lemma}[proposition]{{\bf Lemma}}

\AtBeginDocument{\renewcommand*\contentsname{Contents}}
\AtBeginDocument{\renewcommand*\refname{References}}
%\usepackage{fancyhdr}
\setlength{\headheight}{40pt}
\pagestyle{fancy}
\fancyhead{} % clear all header fields
\fancyhead[L]{
 \begin{tabular}{rl}
    \begin{picture}(25,15)(0,0)
    \put(0,-8){\includegraphics[width=8mm, height=8mm]{hcmut.png}}
    %\put(0,-8){\epsfig{width=10mm,figure=hcmut.eps}}
   \end{picture}&
	%\includegraphics[width=8mm, height=8mm]{hcmut.png} & %
	\begin{tabular}{l}
		\textbf{\bf \ttfamily University of Technology, Ho Chi Minh City}\\
		\textbf{\bf \ttfamily Faculty of Computer Science and Engineering}
	\end{tabular} 	
 \end{tabular}
}
\fancyhead[R]{
	\begin{tabular}{l}
		\tiny \bf \\
		\tiny \bf 
	\end{tabular}  }
\fancyfoot{} % clear all footer fields
\fancyfoot[L]{\scriptsize \ttfamily Assignment for Mathematical Modeling - Academic year 2023 - 2024}
\fancyfoot[R]{\scriptsize \ttfamily Page {\thepage}/\pageref{LastPage}}
\renewcommand{\headrulewidth}{0.3pt}
\renewcommand{\footrulewidth}{0.3pt}


%%%
\setcounter{secnumdepth}{4}
\setcounter{tocdepth}{3}
\makeatletter
\newcounter {subsubsubsection}[subsubsection]
\renewcommand\thesubsubsubsection{\thesubsubsection .\@alph\c@subsubsubsection}
\newcommand\subsubsubsection{\@startsection{subsubsubsection}{4}{\z@}%
                                     {-3.25ex\@plus -1ex \@minus -.2ex}%
                                     {1.5ex \@plus .2ex}%
                                     {\normalfont\normalsize\bfseries}}
\newcommand*\l@subsubsubsection{\@dottedtocline{3}{10.0em}{4.1em}}
\newcommand*{\subsubsubsectionmark}[1]{}
\makeatother


\begin{document}

\begin{titlepage}
\begin{center}
VIETNAM NATIONAL UNIVERSITY, HO CHI MINH CITY \\
UNIVERSITY OF TECHNOLOGY \\
FACULTY OF COMPUTER SCIENCE AND ENGINEERING
\end{center}

\vspace{1cm}

\begin{figure}[h!]
\begin{center}
\includegraphics[width=3cm]{hcmut.png}
\end{center}
\end{figure}

\vspace{1cm}


\begin{center}
\begin{tabular}{c}
\multicolumn{1}{l}{\textbf{{\Large MATHEMATICAL MODELING (CO2011)}}}\\
~~\\
\hline
\\
\multicolumn{1}{l}{\textbf{{Assignment (Semester: 231, Duration: 06 weeks)}}}\\
\\
\textbf{{\Huge "Stochastic Programming}}\\
\\
\textbf{{\Huge and Applications”}}\\[10pt]

\multicolumn{1}{l}{\textbf{{(Version 0.1, in Preparation)}}}\\
\\
\hline
\end{tabular}
\end{center}

\vspace{2cm}

\begin{table}[h]
	\begin{tabular}{rrl}
		\hspace{5 cm} Advisor: & Nguyễn Tiến Thịnh, CSE-HCMUT\\[6pt]
		Students: & 
		\begin{tabular}{@{}ll@{}}
			Trần Đình Đăng Khoa &   2211649 \\
			Trần Đặng Hiển Long &   2252449 \\
			Nguyễn Hồ Phi Ưng &   2252897 \\
			Nguyễn Hồ Đức An &   2252009 \\
			Vũ Minh Quân &   2212828 \\
		\end{tabular}
	\end{tabular}
\end{table}

% \begin{table}[h]
% 	\begin{tabular}{rrl}
% 		\hspace{5 cm} & Advisor: & Nguyễn Tiến Thịnh, CSE-HCMUT\\[6pt]
% 		& Students: & Trần Đình Đăng Khoa - Student 1 ID numbers. \\
% 		& & Trần Đặng Hiển Long - Student 2 ID numbers. \\
% 		& & Nguyễn Hồ Phi Ưng - Student 3 ID numbers. \\
% 		& & Nguyễn Hồ Đức An - 19181716. \\
% 		& & Vũ Minh Quân - 19181716. \\
% 	\end{tabular}
% \end{table}

\vspace*{1cm}
	

\begin{center}
{\footnotesize HO CHI MINH CITY, OCTOBER 2023}
\end{center}
\end{titlepage}


%\thispagestyle{empty}

\newpage
\tableofcontents
\newpage


%%%%%%%%%%%%%%%%%%%%%%%%%%%%%%%%%
\section{Member list \& Workload}

\begin{center}
\begin{tabular}{|c|c|c|l|c|}
\hline
\textbf{No.} & \textbf{Fullname} & \textbf{Student ID} & \textbf{Problems} & \textbf{Percentage of work}\\
\hline 
%%%%%Student 1%%%%%%%%%%
\multirow{3}{*}{1} & \multirow{3}{*}{Trần Đình Đăng Khoa} & \multirow{3}{*}{2211649} & - Relation \& Counting: 1, 2, 3& \multirow{3}{*}{30\%}\\
 & &  & Bonus: 1, 2, 3. &\\
 & &  & - Probability: 1, 2, 3. &\\
\hline 
%%%%%Student 2%%%%%%%%%%%
\multirow{3}{*}{2} & \multirow{3}{*}{Trần Đặng Hiển Long} & \multirow{3}{*}{2252449} & - Relation \& Counting: 4, 5, 6& \multirow{3}{*}{20\%}\\
 & &  & Bonus: 4, 5, 6. &\\
 & &  & - Graph: 1, 2, 3, Bonus: 1, 2, 3. &\\
\hline
%%%%%Student 3%%%%%%%%%%%
\multirow{3}{*}{3} & \multirow{3}{*}{Nguyễn Hồ Phi Ưng} & \multirow{3}{*}{2252897} & - Relation \& Counting: 7, 8, 9& \multirow{3}{*}{20\%}\\
 & &  & Bonus: 7, 8, 9. &\\
 & &  & - Probability: 4, 5, 6. &\\
%%%%%Student 4%%%%%%%%%%%
\hline
\multirow{3}{*}{4} & \multirow{3}{*}{Nguyễn Hồ Đức An} & \multirow{3}{*}{2252009} & - Relation \& Counting: 10, 11, 12& \multirow{3}{*}{20\%}\\
 & &  & Bonus: 10, 11, 12. &\\
 & &  & - Graph: 4, 5, 6, Bonus: 4, 5, 6. &\\
%%%%%Student 5%%%%%%%%%%%
\hline
\multirow{3}{*}{5} & \multirow{3}{*}{Vũ Minh Quân} & \multirow{3}{*}{2212828} & - Relation \& Counting: 13, 14, 15& \multirow{3}{*}{10\%}\\
 & &  & Bonus: 13, 14, 15. &\\
 & &  & - Probability: 7, 8, 9. &\\
\hline
\end{tabular}
\end{center}

%%%%%%%%%%%%%%%%%%%%%%%%%%%%%%%%%
\section{Introduction to Stochastic Programming and Optimization}
	\subsection{Programming? What is Stochastic Programming? Uncertainty?}
	...
	
	\subsection{Basic concepts, assumptions - Motivation}
	...

	\subsection{Bonus exercises}
	...

%%%%%%%%%%%%%%%%%%%%%%%%%%%%%%%%%
\section{Probabilty}
	\subsection{Problem 1}
	...
	
	\subsection{Problem 2}
	...

	\subsection{Bonus exercises}
	...

%%%%%%%%%%%%%%%%%%%%%%%%%%%%%%%%%


\section{Graph}
	\subsection{Problem 1}
	...
	
	\subsection{Problem 2}
	...

	\subsection{Bonus exercises}
	...

\begin{thebibliography}{80}


\bibitem{bib1}
...


\bibitem{bib2}
...


\end{thebibliography}
\end{document}

